\documentclass[10pt,a4paper]{amsart}
\usepackage[margin=19mm]{geometry}
\usepackage{amsmath,amssymb,mathtools}
\usepackage[scr=rsfs,cal=euler]{mathalfa}
\usepackage{xfrac}
\usepackage[unicode,hidelinks,bookmarks=false]{hyperref}
\newcommand{\ie}{i.e.\ }
\newcommand{\cf}{cf.\ }
\newcommand{\resp}{respectively\ }
\newcommand{\Subsection}{\subsection}
\providecommand{\nobreakdash}{\nobreak\hbox{-}}
\setlength{\emergencystretch}{2em}
\multlinegap0pt
\usepackage{bm}
\usepackage{bbm}
\usepackage{dsfont}
\usepackage{xspace}
\usepackage{cancel}
\usepackage{mathtools}
\usepackage{amsthm}
\makeatletter
\@ifundefined{theorem}{\newtheorem{theorem}{Theorem}}{}
\makeatother
\newcommand*\diff{\mathop{}\!\mathrm{d}}
\newcommand{\HH}{\mathsf{H}} 
\newcommand{\II}{\mathsf{I}}
\newcommand{\N}{\mathbb{N}}
\newcommand{\R}{\mathbb{R}} % real numbers
\newcommand{\eqdef}{\stackrel{\mathrm{def}}{=}}
\title[Intrinsic dimensional functional inequalities on model space]{Intrinsic dimensional functional inequalities on model spaces}
\author{Alexandros Eskenazis, Yair Shenfeld}
\date{Source: arXiv:2303.00784v3 (CC BY 4.0)}
\begin{document}
\maketitle

\noindent\emph{Source and licence.} Intrinsic dimensional functional inequalities on model spaces. Version arXiv:2303.00784v3 (CC BY 4.0), distributed under the Creative Commons Attribution 4.0 International licence, as recorded in the arXiv licence metadata for this exact version and shown on the abstract page https://arxiv.org/abs/2303.00784v3. Full rights notice: licenses/arxiv-2303.00784-CC-BY-4.0.txt.\par\smallskip

\noindent\textit{Editorial scope.} Counted unit: the Theorem whose source label is \texttt{thm:cramer-rao} in the article (source0.tex lines 542--554, source label \texttt{thm:cramer-rao}), delivered together with the article's own complete original proof of it (source0.tex lines 558--604, one \texttt{proof} environment of 47 lines). No mathematics is written, repaired or re-derived: statement and proof are the article's own text line for line. The proof is detached from the statement in the source: the article states the result at lines 542--554 and prints its proof at lines 558--604, and the proof environment's own header names the statement, so the association is the article's own. Required context retained uncounted: thm:homogeneous (lines 464--473); eq:regu (lines 525--527). Every definition, display and label of those lines is retained unchanged. Omitted from this package and not redistributed: 1--463, 474--524, 528--541, 555--557, 605--1913. Nothing inside a retained excerpt is omitted or altered: every retained body is delivered exactly as the article's lines are written, so no omission receipt of any kind accompanies this package. Citations: the retained excerpts carry no citation command at all, so no bibliography entry of the article is transcribed and no reference is redistributed. Reference adaptation: none was needed and none was made; every reference inside the delivered unit resolves inside it, so no unresolved reference or citation is produced. Printed slips kept literal: no printed word, symbol, hypothesis or number of the counted statement or of its proof was altered. Layout and class: the article's own class is \texttt{amsart} with packages including fullpage, fontenc, inputenc, hyperref, kpfonts, graphicx, euscript, xcolor, dutchcal; the wrapper is the lane's pinned \texttt{amsart} editorial wrapper at 10pt on A4, which restates the theorem-like environments the retained text needs and adds the article's own macro definitions that the retained text needs (\texttt{\textbackslash HH}, \texttt{\textbackslash II}, \texttt{\textbackslash N}, \texttt{\textbackslash R}, \texttt{\textbackslash diff}, \texttt{\textbackslash eqdef}), with extra wrapper packages (bm, bbm, dsfont, xspace, cancel, mathtools). The delivered numbering is the unit's own. Assets: no retained span contains a figure environment, a graphics-inclusion command, a PDF-page inclusion, a table or a TikZ picture; no image file is copied into this package, the package declares no asset dependency and no image file is redistributed with it. Licence: version arXiv:2303.00784v3 of this article is distributed under the Creative Commons Attribution 4.0 International licence, as recorded in the arXiv licence metadata for this exact version and shown on the abstract page https://arxiv.org/abs/2303.00784v3. A full rights notice is delivered at licenses/arxiv-2303.00784-CC-BY-4.0.txt. Characters: every retained excerpt is pure ASCII, so the delivered engine, XeLaTeX with its default Latin Modern text font, reports no missing character.
\begin{theorem} \label{thm:homogeneous}
Fix $c_1,c_2>0$, $n\in\N$, $p_1,\ldots,p_n\geq0$ and let $\rho$ be a $(p_1,\ldots,p_n)$-homogeneous measure such that for any Borel probability measure $\mu$ on $\R^n$,
\begin{equation} \label{eq:homog-assumption}
\HH(\mu\|\rho)\leq c_1 \II(\mu\|\rho) + c_2.
\end{equation} 
Then, for any Borel probability measure $\mu$ on $\R^n$ with  positive differentiable density $f$, we have
\begin{equation}\label{eq:homog-conclusion}
\HH(\mu\|\rho)\leq  \frac{1}{2} \sum_{k=1}^n (1+p_k)  \log\left( \frac{2ec_1}{1+p_k} \int_{\R^n} \frac{(\partial_k f(y))^2}{f(y)} \,\diff\rho(y)\right) + c_2.
\end{equation}
\end{theorem}

\begin{equation} \label{eq:regu}
\int_\Omega \nabla_\theta f(x;\theta) \,\diff\lambda(x) = 0
\end{equation}

\begin{theorem} \label{thm:cramer-rao}
Fix $c_1,c_2>0$, $n\in\N$, $p_1,\ldots,p_n\geq0$ and let $\rho$ be a $(p_1,\ldots,p_n)$-homogeneous measure  such that for any probability measure $\mu$ on $\R^n$,
\begin{equation} \label{eq:cr-assumption}
\HH(\mu\|\rho)\leq c_1 \II(\mu\|\rho) + c_2.
\end{equation}
Then, for every parametric family $\{\mu_\theta\}_{\theta\in\R^n}$ and every absolutely continuous probability measure $\pi$ on $\R^n$ whose density with respect to $\rho$ is $h:\R^n\to\R_+$, we have
\begin{equation}
\begin{split}
&I\big(\pi;\{\mu_\theta\}\big)+ \HH(\pi\|\rho)
\\& \leq \frac{1}{2}\sum_{k=1}^n (1+p_k) \log\left( \frac{2ec_1}{1+p_k}\Big(\int_{\R^n} \frac{(\partial_{k} h(\theta))^2}{h(\theta)}\,\diff\rho(\theta) + \int_{\R^n} \int_\Omega \frac{(\partial_{\theta_k} f(x;\theta))^2}{f(x;\theta)} \,\diff\lambda(x)\,\diff\pi(\theta)\Big) \right) + c_2.
\end{split}
\end{equation}
\end{theorem}

\begin{proof} [Proof of Theorem \ref{thm:cramer-rao}]
Consider the function $f:\Omega\to\R_+$ given by
\begin{equation} \label{eq:cramer-f}
\forall \ x\in\Omega,\qquad f(x)\eqdef \int_{\R^n} f(x;\theta) \,\diff\pi(\theta)
\end{equation}
and observe that
\begin{equation}
\int_\Omega f(x)\,\diff\lambda(x) = \int_\Omega \int_{\R^n} f(x;\theta) \,\diff\pi(\theta)\,\diff\lambda(x) = \int_{\R^n} \int_\Omega \, \diff\mu_\theta(x) \diff\pi(\theta) = 1.
\end{equation}
Moreover, for $x\in\Omega$, consider the function $h_x:\R^n\to\R_+$ given by
\begin{equation}
\forall \ \theta\in\R^n, \qquad h_x(\theta) \eqdef \frac{h(\theta)f(x;\theta)}{f(x)}
\end{equation}
and notice that the measure $\nu_x$ on $\R^n$ with $\diff\nu_x(\theta) = h_x(\theta)\diff \rho(\theta)$ is a probability measure since
\begin{equation}
\nu_x(\R^n) = \int_{\R^n} \frac{h(\theta) f(x;\theta)}{f(x)} \, \diff\rho(\theta) = \int_{\R^n} \frac{f(x;\theta)}{f(x)} \,\diff\pi(\theta)  \stackrel{\eqref{eq:cramer-f}}{=}1.
\end{equation}
By Theorem \ref{thm:homogeneous} and the assumption on $\rho$, for every $x\in\Omega$ we have
\begin{equation}
\HH(\nu_x\|\rho) \leq \frac{1}{2} \sum_{k=1}^n (1+p_k)\log\Big( \frac{2ec_1}{1+p_k} \int_{\R^n} \frac{(\partial_k h_x(\theta))^2}{h_x(\theta)}\,\diff\rho(\theta)\Big) + c_2.
\end{equation}
Integrating this inequality with respect to the probability measure $f(x)\diff\lambda(x)$, we get
\begin{equation} \label{eq:rao-long}
\begin{split}
\int_\Omega \HH(\nu_x\|\rho)f(x) \,\diff\lambda(x) & \leq \frac{1}{2} \sum_{k=1}^n (1+p_k) \int_\Omega\log\Big( \frac{2ec_1}{1+p_k} \int_{\R^n} \frac{(\partial_k h_x(\theta))^2}{h_x(\theta)}\,\diff\rho(\theta)\Big)f(x)\,\diff\lambda(x) + c_2
\\ & \leq \frac{1}{2} \sum_{k=1}^n (1+p_k) \log\Big( \frac{2ec_1}{1+p_k} \int_\Omega \int_{\R^n} \frac{(\partial_k h_x(\theta))^2}{h_x(\theta)}f(x) \,\diff\rho(\theta)\,\diff\lambda(x)\Big) + c_2,
\end{split}
\end{equation}
where the last line follows from Jensen's inequality. Moreover, by definition we have
\begin{equation}
\begin{split}
\int_\Omega \HH(\nu_x\|\rho)f(x) \,\diff\lambda(x) & = \int_\Omega \int_{\R^n} h(\theta) f(x;\theta) \log\Big( \frac{h(\theta)f(x;\theta)}{f(x)}\Big) \, \diff\rho(\theta)\,\diff\lambda(x) 
\\ & = \int_{\R^n} h(\theta) \log h(\theta)\,\diff\rho(\theta) + \int_{\R^n} \int_\Omega f(x;\theta) \log \big(\frac{f(x;\theta)}{f(x)} \big)\, \diff\lambda(x)\,\diff\pi(\theta)
\\ & = \HH(\pi\|\rho) + I\big(\pi;\{\mu_\theta\}\big).
\end{split}
\end{equation}
Similarly, computing the integral on the right-hand side of \eqref{eq:rao-long}, gives
\begin{equation*}
\begin{split}
 \int_\Omega \int_{\R^n} &\frac{(\partial_k h_x(\theta))^2}{h_x(\theta)}f(x) \,\diff\rho(\theta)\,\diff\lambda(x) =  \int_\Omega \int_{\R^n} \frac{\big(f(x;\theta)\partial_k h(\theta) + h(\theta) \partial_{\theta_k} f(x;\theta)\big)^2}{h(\theta)f(x;\theta)} \,\diff\rho(\theta)\,\diff\lambda(x)
 \\ & = \int_{\R^n}\frac{(\partial_k h(\theta))^2}{h(\theta)} \int_\Omega f(x;\theta) \, \diff\lambda(x)\,\diff\rho(\theta) + 2 \int_{\R^n} \partial_k h(\theta) \int_\Omega \partial_{\theta_k} f(x;\theta) \,\diff\lambda(x) \,\diff\rho(\theta)
 \\ & \qquad\qquad\qquad\qquad\qquad\qquad\qquad\qquad\qquad + \int_{\R^n} h(\theta) \int_\Omega \frac{(\partial_{\theta_k}f(x;\theta))^2}{f(x;\theta)}\,\diff\lambda(x)\,\diff\rho(\theta)
 \\ & \stackrel{\eqref{eq:regu}}{=} \int_{\R^n} \frac{(\partial_{k} h(\theta))^2}{h(\theta)}\,\diff\rho(\theta) + \int_{\R^n} \int_\Omega \frac{(\partial_{\theta_k} f(x;\theta))^2}{f(x;\theta)} \,\diff\lambda(x)\,\diff\pi(\theta).
\end{split}
\end{equation*}
Combining everything, we deduce the desired inequality.
\end{proof}

\end{document}
